% 附录：术语表
\section{术语表}\label{sec:glossary}

下表给出正文反复使用的核心术语，便于非本领域读者快速对齐口径。术语按主题分组，不做穷尽罗列。

\begin{longtable}{@{}>{\raggedright\arraybackslash}p{3.6cm}>{\raggedright\arraybackslash}p{11.6cm}@{}}
\caption{核心术语表}\\
\toprule
术语 & 释义与本文用法 \\
\midrule
\endhead
可交易缝隙 & \textbf{本文工作术语}。指扣除交易成本与冲击成本后、经多重检验校正仍为正的风险调整收益机会。与文献中 mispricing、profitable trading opportunity、limits to arbitrage 相关，但对齐关系见第 \ref{sec:intro} 章的对齐表。 \\
\midrule
定价偏差（mispricing） & 价格与基本面价值（或模型隐含价值）之间的偏离。本文视其为可交易缝隙的原材料，而非其本身。 \\
\midrule
alpha & 相对基准的超额收益。本文强调其“净额”含义：扣成本、扣多重检验校正后仍存在者才算数。 \\
\midrule
信噪比 & 可预测成分与噪声成分的比值。金融预测的信噪比极低（月度截面 $R^2$ 常在 0.5\% 以下），这是过拟合风险的结构性来源。 \\
\midrule
IC / IR & 信息系数/信息比率。IC 为信号与下期收益的截面相关系数，IR 为 IC 的均值与标准差之比。属“信号类”必报指标。 \\
\midrule
DSR（缩水夏普） & 对夏普比率做选择偏差、样本长度与非正态校正后的显著性度量。数值为分位数（越大越显著）。 \\
\midrule
PBO（回测过拟合概率） & 用组合对称交叉验证（CSCV）估计“样本内最优策略在样本外落到中位数以下”的概率。PBO 越高越可疑。 \\
\midrule
walk-forward（滚动前向） & 逐段训练、逐段前向测试的切分方式。与一次性时间切分相对，后者容易被反复用于选模型而失去样本外意义。 \\
\midrule
前视偏差（look-ahead） & 在计算时点使用了当时尚不可得的信息（如用披露日之前的财报、用全样本统计量标准化）。不会报错，只会让结果变好。 \\
\midrule
数据窥探与多重检验 & 测足够多的候选，总能找到“看起来有效”的。校正方法包括 Reality Check、DSR、PBO，以及提高 $t$ 阈值。 \\
\midrule
冲击成本与参与率 & 大额成交对价格的自我冲击。参与率 $=$ 成交额/当日成交量；冲击率常用平方根形式 $k\sqrt{\text{参与率}}$。 \\
\midrule
平方根律 & 市场冲击随参与率近似按平方根增长的经验规律，是本文成本层的默认形态。 \\
\midrule
Kyle $\lambda$ & 微观结构中价格对订单流的敏感度（价格冲击系数）。$\lambda$ 越大，单位信息能换取的利润越薄。 \\
\midrule
信息租金耗散 & 信息被越多参与者使用，其超额收益越薄，机制由 Grossman--Stiglitz 悖论与 Kyle 模型刻画。 \\
\midrule
信息半衰期 & 信号预测力衰减到一半所需的时间。它决定了最优调仓频率（见第 \ref{sec:gapcondition} 节命题）。 \\
\midrule
CVaR & 条件在险价值（尾部损失的期望），用于刻画极端情形下的风险。 \\
\midrule
波动率目标化 & 按预测波动率反向调整仓位，使组合风险保持在目标水平。 \\
\midrule
制度（regime） & 市场所处的宏观/结构状态（如高波动、低波动、趋势、震荡）。制度切换是模型失效的主要来源之一。 \\
\midrule
容量与拥挤度 & 容量指策略能承载的资金规模；拥挤度指多少资金在用相似信号。二者共同决定超额收益的可持续性。 \\
\midrule
LLM 智能体 & 以大型语言模型为核心、能记忆与调用工具并做多步决策的交易代理。本文对其证据强度评价较低（见第 \ref{sec:nlp} 章）。 \\
\midrule
训练语料泄漏 & 用于回测的样本落在模型预训练语料覆盖期内，模型“知道答案”，使回测收益不可信。 \\
\bottomrule
\end{longtable}
